Sobre un metall alcalí incideix llum de longitud d'ona $\lambda = 3,00\times 10^{2}\ \mathrm{nm}$. Si els fotoelectrons emesos tenen una energia cinètica màxima de $2,00\ \mathrm{eV}$, calculeu:

\begin{apartat}{1}
L'energia (en eV) d'un fotó de la llum incident.
\end{apartat}

\begin{apartat}{1}
El treball d'extracció (en eV) corresponent a aquest metall.
\end{apartat}

\dades{%
  $1\ \mathrm{eV} = 1,60\times 10^{-19}\ \mathrm{J}$.\\
  Constant de Planck, $h = 6,63\times 10^{-34}\ \mathrm{J\,s}$.\\
  Velocitat de la llum, $c = 3,00\times 10^{8}\ \mathrm{m\,s^{-1}}$.}
